% Thai-only cover labels; proper names/course title retain their official form.
\def\CoverPublication{การเผยแพร่ข้อสอบ}
\def\CoverAllowed{อนุญาต}\def\CoverNotAllowed{ไม่อนุญาต}
\def\CoverCourse{วิชา}\def\CoverSeat{เลขที่นั่งสอบ}
\def\CoverDateLabel{วันที่สอบ}\def\CoverDate{\ExamDate}
\def\CoverTimeLabel{เวลาสอบ}\def\CoverDuration{\ExamDuration}
\def\CoverGroups{กลุ่มนักศึกษา}\def\CoverRegular{ปกติ}
\def\CoverSet{ชุดข้อสอบ}
\def\CoverName{ชื่อ-นามสกุล :}\def\CoverID{รหัสนักศึกษา :}
\def\CoverDepartment{ภาควิชา/สาขาวิชา :}\def\CoverSection{กลุ่ม :}
\def\CoverInstructionsLabel{คำชี้แจง}\def\CoverConditionsLabel{เงื่อนไขการสอบ}
\def\CoverInstructions{%
1. ข้อสอบมี \textbf{3 ส่วน 12 ข้อ คะแนนเต็ม 100 คะแนน} จำนวน \pageref{LastPage}\ หน้า (รวมหน้าปก)\par
\hspace{.35cm}-- ส่วนที่ 1 ข้อกำหนดและความมั่นคงปลอดภัย ข้อ 1-4 รวม 25 คะแนน\par
\hspace{.35cm}-- ส่วนที่ 2 การวางแผนและดำเนินโครงการ ข้อ 5-7 รวม 25 คะแนน\par
\hspace{.35cm}-- ส่วนที่ 3 แผนภาพ การออกแบบหน้าจอ และระบบควบคุมเวอร์ชัน ข้อ 8-12 รวม 50 คะแนน\par
\hspace{.35cm}\textbf{\underline{ให้ทำคำตอบลงในข้อสอบ}}\par\vspace{.08cm}
2. เขียนชื่อและรหัสนักศึกษาลงในข้อสอบทุกหน้า\par\vspace{.06cm}
3. ใช้ปากกาหมึกสีน้ำเงินหรือสีดำ หากพื้นที่คำตอบไม่พอ ให้เขียนต่อด้านหลังของหน้าข้อสอบข้อเดียวกัน\par\vspace{.06cm}
4. ไม่มีการชี้แจงหรือแก้โจทย์ระหว่างสอบ หากไม่ชัดเจน ให้ตั้งสมมติฐานที่สมเหตุสมผลพร้อมอธิบาย}
\def\CoverConditions{%
1. \textbf{ไม่อนุญาตให้นำเครื่องคำนวณเข้าห้องสอบ}\par\vspace{.08cm}
2. อนุญาตให้นำกระดาษ A4 1 แผ่นเข้าห้องสอบ โดยต้อง\textbf{\underline{เขียนชื่อและรหัสนักศึกษาบนกระดาษ}} และ\textbf{\underline{ส่งพร้อมข้อสอบ}} หากไม่ส่งกระดาษตามเงื่อนไข \textbf{จะหัก 5 คะแนนจากคะแนนสอบ}}
\def\CoverApproval{ข้อสอบได้ผ่านการพิจารณาจากคณะกรรมการประจำภาควิชาวิศวกรรมคอมพิวเตอร์แล้ว}
\def\CoverHead{หัวหน้าภาควิชา}\def\CoverInstructor{อาจารย์ผู้สอน / ผู้ประสานงานรายวิชา}
\def\CoverPhone{เบอร์โทรศัพท์ : ................................................}
\def\CoverInternational{นานาชาติ}
\input{\MockAssetRoot/frontmatter/cover-layout.tex}
